\chapter{Cheat sheet}

\section{Syntax}

\begin{center}
\footnotesize
\begin{tabular}{lp{8.4cm}}
\toprule
\code{;;; text} & comment to end of line\\
\code{value -> x} & assignment; the arrow points at the destination\\
\code{expr =>} & print the result(s), prefixed \code{**}\\
\code{expr ==>} & as \code{=>} but prints structure in full\\
\code{npr(x)} & print \code{x} and a newline\\
\code{a sys\_>< b} & concatenate anything into a string\\
\code{a >< b} & concatenate two strings\\
\code{'text'} & string (mutable, counted, \emph{not} NUL-terminated)\\
\code{\dq word\dq} & word --- an interned symbol\\
\code{[a b c]} & list; contents are \emph{quoted}\\
\code{[\^{}x \^{}(e) \^{}\^{}xs]} & \ldots with \code{x}, \code{e} evaluated and \code{xs} spliced\\
\code{\{1 2 3\}} & vector\\
\code{[\% e \%]} \ \code{\{\% e \%\}} & list/vector of everything \code{e} leaves on the stack\\
\code{p(\% a \%)} & closure: \code{p} partially applied to \code{a}\\
\bottomrule
\end{tabular}
\end{center}

\section{Blocks --- every opener has a closer}

\begin{center}
\footnotesize
\begin{tabular}{ll}
\toprule
\code{define f(a, b); \ldots enddefine;} & \code{if \ldots then \ldots elseif \ldots else \ldots endif}\\
\code{define f(a) -> r; \ldots enddefine;} & \code{unless \ldots then \ldots endunless}\\
\code{procedure(a); \ldots endprocedure} & \code{for i from 1 to n do \ldots endfor}\\
\code{define :class C; slot s = v; enddefine;} & \code{for x in list do \ldots endfor}\\
\code{define :method m(x:C); \ldots enddefine;} & \code{while c do \ldots endwhile}\\
\code{section \$-name => exports; \ldots endsection;} & \code{repeat \ldots quitif(c); \ldots endrepeat}\\
\code{exload tag [libs] \ldots endexload;} & \code{switchon e case < 0 then \ldots endswitchon}\\
\bottomrule
\end{tabular}
\end{center}

\section{Declarations}

\code{vars} permanent (global) --- required for matcher pattern variables.
\code{lvars} lexical local --- the default choice inside a procedure.
\code{constant}, \code{lconstant} their immutable forms.
\code{dlocal x = v} rebind for the dynamic extent of the call, restored on
\emph{any} exit.

\section{Mishap decoder}

\begin{center}
\footnotesize
\begin{tabular}{lp{7.9cm}}
\toprule
\code{DECLARING VARIABLE x} & (a warning) you used an undefined name ---
  usually a typo or a missing \code{define}\\
\code{NUMBER(S) NEEDED} & arithmetic on a non-number, often the value of
  that undefined name\\
\code{STRING NEEDED} & a word \code{\dq x\dq} was passed where a string
  \code{'x'} was wanted --- check the quotes\\
\code{ILLEGAL DECLARATION OF WORD} & the name is already a protected
  system identifier; rename\\
\code{MISPLACED EXPRESSION ITEM} & syntax: usually a missing keyword or a
  construct used with the wrong shape\\
\code{EXECUTING NON-PROCEDURE} & calling something that is not a procedure
  --- often an undefined name\\
\code{CAN'T RESTORE - NOT SAME SYSTEM AND VERSION} & the \code{.psv} image
  predates the current \code{basepop11}; rebuild the images\\
\code{sysEXECUTE: NOT AT EXECUTE LEVEL} & VM planting attempted at
  compile-time top level --- wrap it in a procedure\\
\bottomrule
\end{tabular}
\end{center}

\section{Files, processes, tables}

\begin{pop}
;;; read a file line by line
vars dev = sysopen('/path/file', 0, "line");
vars rep = line_repeater(dev, inits(4096));   ;;; buffer size = max line length
repeat
    rep() -> line;
    quitif(line == termin);
    if issubstring('ERROR', 1, line) then ... endif;
endrepeat;

sysobey('ls /tmp > /tmp/out');                          ;;; run a command
vars r = sys_obey_linerep('set +e; grep pat file');     ;;; ...and stream it

newproperty([], 50, false, true) -> tbl;   ;;; hash table, keys matched by ==
newmapping([],  50, false, true) -> m;     ;;; ...keys matched by =
\end{pop}

\code{line\_repeater} truncates lines longer than its buffer, so size
\code{inits(n)} generously. \code{sys\_obey\_linerep} runs its command line
under \code{errexit}: prefix \code{set +e;} whenever a failure is expected.

\section{Regular expressions}

The escape character is \code{@}; \code{.\ * [ ] \^{} \$} are literal unless
escaped; there is no alternation and no lookaround.

\begin{pop}
vars err, p;
regexp_compile('ERROR @[0-9@]@{1,@}') -> (err, p);   ;;; test err -- false = ok
p(1, 'ERROR 42 happened', false, false) -> (i, n);   ;;; -> 1  8
substring(i, n, 'ERROR 42 happened') =>              ;;; ** ERROR 42
\end{pop}

Case-insensitive: \code{regexp\_compile(s, 1, false)}, or \code{@i} in the
pattern. \code{p} is a compiled procedure --- define it once, reuse it.

\section{Running things}

\begin{center}
\footnotesize
\begin{tabular}{ll}
\toprule
\code{./poplog pop11} & Pop-11 REPL\\
\code{./poplog basepop11 -target/psv/prolog.psv} & Prolog\\
\code{./poplog basepop11 -target/psv/clisp.psv} & Common Lisp\\
\code{./poplog basepop11 +target/psv/pml.psv} & Standard ML\\
\code{tools/forth.sh} & Forth\\
\code{./poplog basepop11 prog.p} & run a Pop-11 file\\
\code{tools/pop11-lsp} \ \code{tools/pop11-swank} \ \code{tools/pop11-mcp} & editor and agent servers\\
\bottomrule
\end{tabular}
\end{center}
